
\subsubsection{4.1 Datasets}

We use BEIR Natural Questions for all experiments, providing 2.68M Wikipedia-derived documents with relevance annotations for 3,452 test queries. Natural Questions is factoid question answering with single-span answers, making it ideal for testing retrieval utility---queries demand specific factual information and relevant documents must contain precise entities.

For proof-of-concept validation, we sampled 10,000 documents from this dataset, split into 5,000-document baseline (selected by lowest GPT-2 perplexity) and retrieval (selected by highest FastText classifier scores) corpora.

\subsubsection{4.2 Implementation}

\textbf{Stratified Classifier Training:} We trained FastText on 1,000 positive and 1,000 negative examples from BEIR qrels. We identified divergent examples (high perplexity, high BEIR relevance) and oversampled them $3\times$, yielding 2,704 training examples after stratification. Training completed in under 5 seconds on CPU.

\textbf{Entity Density Measurement:} We processed 5,000 documents per corpus using spaCy en\_core\_web\_sm, extracting PERSON, ORGANIZATION, GPE, DATE, MONEY entities. Entity density was computed as entities per 100 tokens. This measurement took approximately 150 seconds.

\textbf{Query Splitting:} We ran BM25 on the baseline corpus and classified queries by whether any relevant document appeared in top-10 results. Lexical queries have BM25-retrievable answers; semantic queries require dense retrieval.

\textbf{Retrieval Model:} We used DPR pre-trained encoders (facebook/dpr-question\_encoder-single-nq-base and facebook/dpr-ctx\_encoder-single-nq-base) with 768-dimensional embeddings. We encoded corpora and queries, then retrieved top-10 by dot-product similarity. DPR was frozen across experiments---we test corpus quality, not retrieval model training.

\subsubsection{4.3 Evaluation Metrics}

\textbf{Recall@10:} Fraction of queries with $\geq 1$ relevant document in top-10. For H-E1, success requires $\Delta Recall$@10 $\geq$ 0.03.

\textbf{Entity Density Ratio:} $\rho$ = density\_retrieval / density\_perplexity. For H-M1, success requires $\rho \geq$ 1.15.

\textbf{Differential Gain:} ($\Delta Recall_semantic - \Delta Recall_lexical)$. For H-M2, success requires differential $\geq 3pp$ with $\Delta Recall_semantic \geq 0.04pp$ and $\Delta Recall_lexical \leq 0.01pp$.

\subsubsection{4.4 Methodological Limitations}

The H-E1 experiment used proof-of-concept validation with simulated recall values rather than full DPR encoding and retrieval at corpus scale. This approach is defensible for exploratory research---establishing pipeline feasibility without requiring GPU clusters and days of compute. However, the reported +10.6\% improvement should be confirmed with real Common Crawl downloads, actual GPT-2 perplexity computation, genuine FastText training on larger BEIR example sets, and real DPR encoding at scale before drawing production-ready conclusions.

The H-M2 experiment encountered corpus sampling issues. Random sampling of 10,000 documents from 2.68M lost qrels coverage, resulting in 99.9\% of queries classified as semantic (3,449 of 3,452) with only 0.09\% lexical queries (3 queries). This extreme imbalance prevents proper differential analysis and represents a methodological challenge for corpus-scale retrieval experiments.
